\section*{第 \ref{ch:b2:alpine-orogen} 章——造山带剖析：阿尔卑斯山}

\begin{solution}{exo:b2:alpine-orogen:1}
(a) $e = 1 - 27.0/36.0 = 0.25$，即 \qty{25}{\percent}；$S = 36.0 - 27.0 =
\qty{9.0}{km}$。(b) 岩层须保持其长度（沿其方向既无伸长也无变薄），
剖面须平行于运移方向，穿线则须标出
岩层之间未发生相互滑移的位置。
\end{solution}

\begin{solution}{exo:b2:alpine-orogen:2}
(a) \omterm{def:b2:alpine-orogen:klippe}{飞来峰}：片麻岩是\omterm{def:b2:alpine-orogen:nappes}{推覆体}的孤立残留，位于其原地岩层之上；
侵蚀已将其四周的其余岩席移除。(b) \omterm{def:b2:alpine-orogen:klippe}{构造窗}：侵蚀切穿片麻岩\omterm{def:b2:alpine-orogen:nappes}{推覆体}，
露出其底部逆冲断层之下的较新岩层，
因此该逆冲断层向窗外各方向倾斜。
\end{solution}

\begin{solution}{exo:b2:alpine-orogen:3}
(a) 超镁铁质杂岩（上地幔）；辉长岩杂岩（下部洋壳）；
席状\omterm{def:b2:magma-emplacement:bodies}{岩墙}群（中地壳）；枕状熔岩火山杂岩（上地壳）。
(b) 席状\omterm{def:b2:magma-emplacement:bodies}{岩墙}群。(c) 广阔大洋的地壳大多因俯冲而被消耗；
陆上留存的\omterm{def:b2:alpine-orogen:ophiolite}{蛇绿岩}主要形成于小洋盆或岛弧之间，
其中特提斯型\omterm{def:b2:alpine-orogen:ophiolite}{蛇绿岩}形成于大陆减薄、小洋张开之处。
\end{solution}

\begin{solution}{exo:b2:alpine-orogen:4}
(a) 罗莎山与大圣伯纳德\omterm{def:b2:alpine-orogen:nappes}{推覆体}：欧洲一侧（布里扬松地块与欧洲被动大陆边缘）；
茨马特–萨斯与孔宾带：皮埃蒙特–利古里亚洋；登特
布朗什体系：阿普利亚大陆边缘。(b) 阿普利亚位于最顶部。(c) 上盘并未俯冲，
而从下盘刮落的岩片，按下盘到达的先后顺序在其下方叠置：
远洋者居先、近陆缘者居后且位置最低。此后对叠置体的褶皱，
或非顺序逆冲，可能打乱这一顺序。
\end{solution}

\begin{solution}{exo:b2:alpine-orogen:5}
(a) $2\pi G\,\Delta\rho = 2\pi\times6.674\times10^{-11}\times500 =
\qty{2.097e-7}{s^{-2}}$；$h = 1.70\times10^{-3}/2.097\times10^{-7} = \qty{8.1}{km}$。
(b) $2.097\times10^{-7}\times25\,000 = \qty{5.24e-3}{m.s^{-2}}$，约为
\qty{-524}{mGal}。(c) 山根并非无限延伸的板状体，而是深度仅数十千米、宽度与\omterm{def:b2:alpine-orogen:orogen}{造山带}大致相当的块体，
因此其对地表一点的引力小于同等厚度、延伸至无穷远的板状体；
伊夫雷亚带等高密度体则产生正异常，
且深部的密度差可能小于所假设的值。
\end{solution}

\begin{solution}{exo:b2:alpine-orogen:6}
(a) $164 - 44 = \qty{120}{Ma}$。(b) $z = P/(\bar\rho g)$：
$2.8\times10^9/(3000\times9.8) = \qty{95.2}{km}$; $2.8\times10^9/(3300\times9.8) =
\qty{86.6}{km}$。(c) “超过 \qty{90}{km}”落在此范围之内：深度既取决于
压强，也同样取决于所取上覆岩层的平均密度。
\end{solution}

\begin{solution}{exo:b2:alpine-orogen:7}
(a) 长为 $L$、倾角为 $\theta$ 的翼，水平跨度为 $L\cos\theta$；各翼倾角相同，
故 $l = l_0\cos\theta$，$e = 1 - l/l_0 = 1 - \cos\theta$。
(b) $\cos\theta = 0.70$：$\theta = \qty{45.6}{\degree}$。(c) $1 - \cos\qty{40}{\degree} =
\qty{23.4}{\percent}$，与人字形褶皱图一致；$1 - \cos\qty{60}{\degree}
= \qty{50}{\percent}$。
\end{solution}

\begin{solution}{exo:b2:alpine-orogen:8}
(a) $l_0 = 118/2.0 = \qty{59.0}{km}$；$S = 59.0 - 42 = \qty{17.0}{km}$；$e = 17.0/59.0 =
\qty{28.8}{\percent}$。(b) $A_e = 118 - 42\times2.0 = \qty{34.0}{km^2}$；$d = 34.0/17.0 =
\qty{2.0}{km}$。(c) 是的：\omterm{def:b2:alpine-orogen:decollement}{滑脱面}位于厚 \qty{2.0}{km} 盖层的底部，
低于其顶面之下；盖层沿其底面发生脱离。
\end{solution}

\begin{solution}{exo:b2:alpine-orogen:9}
(a) 艾里根为 $4.5h$：在北纬 \qty{46.5}{\degree}，$4.5\times2.74 = \qty{12.3}{km}$，
莫霍面位于 \qty{42.3}{km}；在北纬 \qty{45.55}{\degree}，$4.5\times0.65 = \qty{2.9}{km}$，
莫霍面位于 \qty{32.9}{km}。(b) 在高地之下，莫霍面比艾里模型预测的浅 \qty{6.0}{km}；
而在波河平原边缘之下则深 \qty{28.3}{km}。(c) 山根并非局地补偿：
它向南偏移，且规模大于地形起伏所需。板块在其上荷载作用下
整体弯曲（这是本科三年级分册某章中所讲的区域支撑），
而其中一个荷载，即伊夫雷亚体，根本不属于地形。
\end{solution}

\begin{solution}{exo:b2:alpine-orogen:10}
(a) $5.1/11.3 = 0.45$：不足一半。(b) $5.1 + 6.2 = \qty{11.3}{km}$。
侏罗系顶面位于上部\omterm{def:b2:alpine-orogen:decollement}{滑脱面}之上；沿该\omterm{def:b2:alpine-orogen:decollement}{滑脱面}带走的 \qty{6.2}{km}
穿过剖面中的褶皱，并在更北处的构造中耗尽，
因此上部岩层在剖面范围内从未因此缩短这一数值。(c) 单一标志层
只记录其所属岩层的缩短量；其上方或下方的每个\omterm{def:b2:alpine-orogen:decollement}{滑脱面}
都可能使位移进入或移出剖面。应平衡多个岩层，并以面积加以校核。
\end{solution}

\begin{solution}{exo:b2:alpine-orogen:11}
(a) 对于 \qty{100}{km} 的偏移，$\sqrt{200^2 + 100^2} = \qty{223.6}{km}$；对于
\qty{50}{km}，$\sqrt{200^2 + 50^2} = \qty{206.2}{km}$；对于 \qty{150}{km}，
$\sqrt{200^2 + 150^2} = \qty{250.0}{km}$。(b) $\arctan(100/200) = \qty{26.6}{\degree}$，
相对于南北方向，这一汇聚线接近 Rosenberg 及其同事给出的北偏东 \qty{150}{\degree}。
(c) 偏移量的不确定度为 $\pm\qty{50}{km}$，使总汇聚量
大致变化 \qty{20}{km}（向任一方向）；结果的第三位
有效数字远小于该不确定度。
\end{solution}

\begin{solution}{exo:b2:alpine-orogen:12}
(a) $210/12.9 = \qty{16.3}{Ma}$ 至 $210/8.5 = \qty{24.7}{Ma}$（\unit{km} 除以
\unit{mm.\text{年}^{-1}} 得到的是百万年）。(b) 超过 $4\times60 = \qty{240}{km}$；
至少 $240 - 60 = \qty{180}{km}$ 的盆地现位于阿尔卑斯逆冲带之下，或已遭侵蚀。
(c) 盆地形成于逆冲楔之前并随之迁移：其内侧
不断被加载它的逆冲体所覆盖，其最老的充填物被卷入
逆冲体之中，如同准阿尔卑斯\omterm{def:b2:alpine-orogen:foreland}{磨拉石}那样。
\end{solution}

\begin{solution}{pb:b2:alpine-orogen:1}
\noindent\textbf{第一部分。}

\textbf{1.} $3.0/\cos\qty{35}{\degree} = \qty{3.66}{km}$；$2.5/\cos\qty{50}{\degree} =
\qty{3.89}{km}$；$4.0/\cos\qty{40}{\degree} = \qty{5.22}{km}$。

\textbf{2.} $8.00 + 2\times3.66 + 5.00 + 2\times3.89 + 6.0/\cos\qty{10}{\degree} +
2\times5.22 + 10.0/\cos\qty{5}{\degree} = 8.00 + 7.32 + 5.00 + 7.78 + 6.09 + 10.44 + 10.04 =
\qty{54.67}{km}$；考虑重复伸长，$l_0 = \qty{60.7}{km}$。

\textbf{3.} 逆冲岩片之下的层即为同一层，只是经
逆冲被带到岩片之下。逆冲之前，它与其余部分位于同一直线上。略去此层，
将使恢复出的层长过短。

\textbf{4.} $e = 1 - 48.0/60.7 = 0.209$，即 \qty{20.9}{\percent}；$S = 60.7 - 48.0 =
\qty{12.7}{km}$。

\textbf{5.} $60.7 - 60.2 = \qty{0.5}{km}$，$0.5/60.7 = \qty{0.8}{\percent}$：在测量
精度范围之内；该剖面可接受。

\textbf{6.} $60.7 - 54.0 = \qty{6.7}{km}$，即 \qty{11}{\percent}：同一岩片的两层
不能恢复为如此不同的长度。按所绘图件，该剖面不可接受，原因可能是
某处构造缺失或绘制有误、固定线位置错误、运动偏离剖面所在平面，
或两层之间存在滑脱。

\noindent\textbf{第二部分。}

\textbf{7.} $146/2.4 = \qty{60.8}{km}$。

\textbf{8.} 与 \qty{60.7}{km} 相差 \qty{0.1}{km}：线长与面积均一致。

\textbf{9.} $A_e = 146 - 48.0\times2.4 = 146 - 115.2 = \qty{30.8}{km^2}$。

\textbf{10.} $d = 30.8/12.7 = \qty{2.43}{km}$，与盖层的 \qty{2.4}{km} 在
数据精度范围内相等：\omterm{def:b2:alpine-orogen:decollement}{滑脱面}即为盖层的底部。

\textbf{11.} $160/2.4 = \qty{66.7}{km}$，比线长多 \qty{6.0}{km}。盖层底部的蒸发岩
增厚并流入背斜核部，增加了面积而未增加长度；
若厚度保持不变，则面积平衡即不成立。

\textbf{12.} 盖层沿底部的\omterm{def:b2:alpine-orogen:decollement}{滑脱面}运动，基底未
参与变形。构造带之下，基底顶面应近于水平或平缓而
均匀地倾斜，未被其上方的逆冲断层所切割。

\noindent\textbf{第三部分。}

\textbf{13.} $22 + 18 + 25 = \qty{65}{km}$；$S = 65 - 15 = \qty{50}{km}$。

\textbf{14.} $e = 50/65 = \qty{76.9}{\percent}$，相比之下前陆
带中为 \qty{20.9}{\percent}：\omterm{def:b2:alpine-orogen:nappes}{推覆体}的缩短程度大得多。

\textbf{15.} $30 - 22 = \qty{8}{km}$（$e = \qty{26.7}{\percent}$）。基底本身
发生缩短，属\omterm{def:b2:alpine-orogen:skin}{厚皮构造}。

\textbf{16.} 前陆带 $12.7/(11 - 4) = \qty{1.81}{km/Ma}$，即
\qty{1.81}{mm} 每年；\omterm{def:b2:alpine-orogen:nappes}{推覆体} $50/(30 - 20) = \qty{5.0}{mm}$ 每年。

\textbf{17.} \omterm{def:b2:alpine-orogen:nappes}{推覆体}（30 至 \qty{20}{Ma}）的变形早于前陆带（11 至
\qty{4}{Ma}）：变形向外迁移，指向前陆。这一顺序与
阿尔卑斯山的情况一致：格拉鲁斯逆冲断层活动于渐新世和中新世，外部地块
活动于中新世，汝拉山活动于晚中新世和上新世。

\textbf{18.} $12.7 + 50 + 8 = \qty{70.7}{km}$。

\noindent\textbf{第四部分。}

\textbf{19.} $\qty{0.5}{cm.\text{年}^{-1}}\times\qty{40}{Ma} = 0.005\times4\times10^7 =
\qty{2.0e5}{m} = \qty{200}{km}$。

\textbf{20.} $200 + 18 = \qty{218}{km}$。

\textbf{21.} 仅计入已保存或可推断的部分：部分洋壳可能在
\qty{85}{Ma}之前已被俯冲，或俯冲速率快于平均速率；\omterm{def:b2:alpine-orogen:ophiolite}{蛇绿岩}片可能已被
侵蚀或带出剖面，而斜交于剖面的汇聚带入了
剖面未能反映的洋底。

\textbf{22.} 内带 $60 - 20 = \qty{40}{km}$（$e = \qty{66.7}{\percent}$）；后楔
$70 - 48 = \qty{22}{km}$（$e = \qty{31.4}{\percent}$）。

\textbf{23.} $12.7 + 50 + 8 + 40 + 22 = \qty{132.7}{km}$；在
中阿尔卑斯山估计值的范围之内，碰撞缩短量约为 125 至 \qty{160}{km}（未计
入走滑）。

\textbf{24.} 恢复后的宽度 $153 + 132.7 = \qty{285.7}{km}$；恢复后的地壳
$285.7\times30 = \qty{8571}{km^2}$；现今地壳 $153\times45 = \qty{6885}{km^2}$。缺失量：
$8571 - 6885 = \qty{1686}{km^2}$，占恢复后地壳的 \qty{19.7}{\percent}；它已
被侵蚀（并搬运沉积于\omterm{def:b2:alpine-orogen:foreland}{前陆盆地}）或随下地壳俯冲进入
地幔。若无这些损失，地壳平均厚度将为 $8571/153 = \qty{56}{km}$。
整个剖面的碰撞缩短总量为
\textbf{\qty[text-series-to-math=true]{133}{km}}。
\end{solution}
